%%% accessibility-test.tex
%%% Example document demonstrating accessible graphics with alt text
%%%
%%% IMPORTANT: This document requires:
%%%   - TeX Live 2025 or newer
%%%   - LuaLaTeX compiler (recommended) or pdfLaTeX
%%%   - Biber for bibliography processing
%%%
%%% Compile with:
%%%   lualatex accessibility-test.tex
%%%   biber accessibility-test
%%%   lualatex accessibility-test.tex
%%%   lualatex accessibility-test.tex

%% ============================================================================
%% CRITICAL: DocumentMetadata MUST come before \documentclass
%% This enables PDF tagging for accessibility
%% ============================================================================
%% Note: Use testphase for TeX Live 2025 (Jan), tagging=on for later releases
\DocumentMetadata{
  lang=en-US,
  pdfversion=2.0,
  testphase={phase-III}
}

\documentclass[12pt,letterpaper]{article}

%% ============================================================================
%% Packages
%% ============================================================================
\usepackage[utf8]{inputenc}
\usepackage[T1]{fontenc}
\usepackage{graphicx}
\usepackage{hyperref}
\usepackage{csquotes}
\usepackage{etoolbox}

% BibLaTeX setup
\usepackage[
  backend=biber,
  giveninits=true,
  maxnames=100,
  maxcitenames=5,
  sorting=nyt,
  backref=true
]{biblatex}

% Load our test metadata
\addbibresource{test-metadata.bib}
% Load figures-source metadata for \includegraphicsource test
\addbibresource{../book/optimization/figures/figures-source/00_METADATA.bib}
% Load figures-static metadata for \includegraphicstatic test
\addbibresource{../book/optimization/figures/figures-static/00_METADATA.bib}

% Graphics path - include both direct paths and paths used by macros
\graphicspath{{../optimization/figures/figures-static/}{../optimization/figures/figures-source/}{../}}

%% ============================================================================
%% Load the accessible graphics macros
%% ============================================================================
%%% preamble-accessible-graphics.tex
%%% Accessible graphics commands with alt text support
%%% Uses modern LaTeX tagging (TeX Live 2025+) and BibLaTeX metadata
%%%
%%% Alt text is pulled from the 'abstract' field in the .bib metadata files.
%%%
%%% Requirements:
%%%   - \DocumentMetadata{tagging=on} before \documentclass
%%%   - BibLaTeX with biber backend
%%%   - TeX Live 2025 or newer (for alt= key support)

\usepackage{xargs}
\usepackage{xstring}

%% ============================================================================
%% Helper command to get alt text from bib entry
%% Uses the 'abstract' field which stores our alt text descriptions
%% ============================================================================

% Command to extract alt text using BibLaTeX's internal mechanisms
\makeatletter

% Storage macro for the extracted alt text
\newcommand{\OO@alttext}{}

% Create a format that stores instead of prints
\DeclareFieldFormat{OOalttext}{%
  \gdef\OO@alttext{#1}%
}

% Command to extract alt text from a bib entry
% Uses a custom cite command to access the abstract field
\DeclareCiteCommand{\OOgetabstract}
  {}
  {\printfield[OOalttext]{abstract}}
  {}
  {}

% Helper to include graphics with expanded alt text
% This uses \expanded to force expansion of \OO@alttext before passing to \includegraphics
\newcommand{\OO@includegraphicswithalt}[3]{%
  % #1 = options, #2 = alt text macro, #3 = file path
  \expanded{\noexpand\includegraphics[#1, alt={#2}]{#3}}%
}

\makeatother

%% ============================================================================
%% Footnote metadata command (same as original)
%% ============================================================================

\DeclareRobustCommand\footnotetextmetadata[1]{%
  \footnotetext{\citetitle*{#1}, from \citeurl{#1}.
    \citeauthor{#1}, \citedate{#1}.
  }
}

%% ============================================================================
%% ACCESSIBLE INLINE GRAPHICS (no figure environment)
%% ============================================================================

% \includegraphicstatic[options]{bibkey}
% Includes image from figures-static/ folder with alt text from .bib abstract field
\makeatletter
\newcommandx{\includegraphicstatic}[2][1={width=.4\linewidth}, usedefault]{%
  \OOgetabstract{#2}%
  \begin{center}
  \OO@includegraphicswithalt{#1}{\OO@alttext}{optimization/figures/figures-static/#2}\\
  \end{center}
  \centerline{\scriptsize \copyright~\citeauthor{#2}\protect\footnotemark}
  \footnotetextmetadata{#2}%
}

% \includegraphicsource[options]{bibkey}
% Includes image from figures-source/ folder with alt text from .bib abstract field
\newcommandx{\includegraphicsource}[2][1={width=.4\linewidth}, usedefault]{%
  \OOgetabstract{#2}%
  \begin{center}
  \OO@includegraphicswithalt{#1}{\OO@alttext}{optimization/figures/figures-source/#2}\\
  \end{center}
  \centerline{\scriptsize \copyright~\citeauthor{#2}\protect\footnotemark}
  \footnotetextmetadata{#2}%
}
\makeatother

%% ============================================================================
%% NO-CITE VERSIONS (no copyright line or footnote metadata)
%% Use these for your own graphics where attribution isn't needed
%% ============================================================================

\makeatletter
% \includegraphicstaticnocite[options]{bibkey}
% Includes image from figures-static/ with alt text but no citation/metadata
\newcommandx{\includegraphicstaticnocite}[2][1={width=.4\linewidth}, usedefault]{%
  \OOgetabstract{#2}%
  \begin{center}
  \OO@includegraphicswithalt{#1}{\OO@alttext}{optimization/figures/figures-static/#2}%
  \end{center}
}

% \includegraphicsourcenocite[options]{bibkey}
% Includes image from figures-source/ with alt text but no citation/metadata
\newcommandx{\includegraphicsourcenocite}[2][1={width=.4\linewidth}, usedefault]{%
  \OOgetabstract{#2}%
  \begin{center}
  \OO@includegraphicswithalt{#1}{\OO@alttext}{optimization/figures/figures-source/#2}%
  \end{center}
}
\makeatother

%% ============================================================================
%% ACCESSIBLE FIGURE ENVIRONMENTS (with caption and label)
%% ============================================================================

% \includefigurestatic[caption][options][placement]{bibkey}
% Full figure with caption (from argument or bib title), label, and alt text
\makeatletter
\newcommandx{\includefigurestatic}[4][1=\relax,2={width=.4\linewidth},3=t, usedefault]{%
  \OOgetabstract{#4}%
  \begin{figure}[#3]%
    \centering%
    \OO@includegraphicswithalt{#2}{\OO@alttext}{optimization/figures/figures-static/#4}\\
    \centerline{\scriptsize \copyright~\citeauthor{#4}\protect\footnotemark}
    \caption{%
      \ifx#1\relax % Caption from TITLE field in .bib
      \citetitle*{#4}%
      \else        % Caption provided as an argument to the macro
      #1%
      \fi%
    }%
    \label{fig:#4}%
  \end{figure}%
  \footnotetextmetadata{#4}%
}

% \includefiguresource[caption][options][placement]{bibkey}
% Full figure from figures-source/ folder
\newcommandx{\includefiguresource}[4][1=\relax,2={width=.4\linewidth},3=t, usedefault]{%
  \OOgetabstract{#4}%
  \begin{figure}[#3]%
    \centering%
    \OO@includegraphicswithalt{#2}{\OO@alttext}{optimization/figures/figures-source/#4}\\
    \centerline{\scriptsize \copyright~\citeauthor{#4}\protect\footnotemark}
    \caption{%
      \ifx#1\relax % Caption from TITLE field in .bib
      \citetitle*{#4}%
      \else        % Caption provided as an argument to the macro
      #1%
      \fi%
    }%
    \label{fig:#4}%
  \end{figure}%
  \footnotetextmetadata{#4}%
}

%% --- NO-CITE FIGURE ENVIRONMENTS ---
%% These create figures with captions but no copyright line or footnote metadata

% \includefigurestaticnocite[caption][options][placement]{bibkey}
% Figure with caption but no citation/metadata - for your own graphics
\newcommandx{\includefigurestaticnocite}[4][1=\relax,2={width=.4\linewidth},3=t, usedefault]{%
  \OOgetabstract{#4}%
  \begin{figure}[#3]%
    \centering%
    \OO@includegraphicswithalt{#2}{\OO@alttext}{optimization/figures/figures-static/#4}%
    \caption{%
      \ifx#1\relax % Caption from TITLE field in .bib
      \citetitle*{#4}%
      \else        % Caption provided as an argument to the macro
      #1%
      \fi%
    }%
    \label{fig:#4}%
  \end{figure}%
}

% \includefiguresourcenocite[caption][options][placement]{bibkey}
% Figure from figures-source/ with caption but no citation/metadata
\newcommandx{\includefiguresourcenocite}[4][1=\relax,2={width=.4\linewidth},3=t, usedefault]{%
  \OOgetabstract{#4}%
  \begin{figure}[#3]%
    \centering%
    \OO@includegraphicswithalt{#2}{\OO@alttext}{optimization/figures/figures-source/#4}%
    \caption{%
      \ifx#1\relax % Caption from TITLE field in .bib
      \citetitle*{#4}%
      \else        % Caption provided as an argument to the macro
      #1%
      \fi%
    }%
    \label{fig:#4}%
  \end{figure}%
}
\makeatother

%% ============================================================================
%% ACCESSIBLE FIGURE WITH AUTO-REFERENCE
%% ============================================================================

% \refincludefigurestatic[caption][options][placement]{bibkey}
% Same as includefigurestatic but also auto-references the figure
\makeatletter
\newcommandx{\refincludefigurestatic}[4][1=\relax,2={width=.4\linewidth},3=t, usedefault]{%
  \OOgetabstract{#4}%
  \begin{figure}[#3]%
    \centering%
    \OO@includegraphicswithalt{#2}{\OO@alttext}{optimization/figures/figures-static/#4}\\
    \centerline{\scriptsize \copyright~\citeauthor{#4}\protect\footnotemark}
    \caption{%
      \ifx#1\relax
      \citetitle*{#4}%
      \else
      #1%
      \fi%
    }%
    \label{fig:#4}%
  \end{figure}%
  \autoref{fig:#4}\footnotetextmetadata{#4}%
}

% \refincludefiguresource[caption][options][placement]{bibkey}
% Same as includefiguresource but also auto-references the figure
\newcommandx{\refincludefiguresource}[4][1=\relax,2={width=.4\linewidth},3=t, usedefault]{%
  \OOgetabstract{#4}%
  \begin{figure}[#3]%
    \centering%
    \OO@includegraphicswithalt{#2}{\OO@alttext}{optimization/figures/figures-source/#4}\\
    \centerline{\scriptsize \copyright~\citeauthor{#4}\protect\footnotemark}
    \caption{%
      \ifx#1\relax
      \citetitle*{#4}%
      \else
      #1%
      \fi%
    }%
    \label{fig:#4}%
  \end{figure}%
  \autoref{fig:#4}\footnotetextmetadata{#4}%
}
\makeatother

%% ============================================================================
%% MANUAL ALT TEXT OVERRIDE VERSIONS
%% For cases where you want to specify alt text directly instead of from .bib
%% ============================================================================

% \includegraphicstaticalt[options]{bibkey}{alt text}
% Same as includegraphicstatic but with manual alt text
\newcommandx{\includegraphicstaticalt}[3][1={width=.4\linewidth}, usedefault]{%
  \begin{center}
  \includegraphics[#1, alt={#3}]{optimization/figures/figures-static/#2}\\
  \end{center}
  \centerline{\scriptsize \copyright~\citeauthor{#2}\protect\footnotemark}
  \footnotetextmetadata{#2}%
}

% \includegraphicsourcealt[options]{bibkey}{alt text}
% Same as includegraphicsource but with manual alt text
\newcommandx{\includegraphicsourcealt}[3][1={width=.4\linewidth}, usedefault]{%
  \begin{center}
  \includegraphics[#1, alt={#3}]{optimization/figures/figures-source/#2}\\
  \end{center}
  \centerline{\scriptsize \copyright~\citeauthor{#2}\protect\footnotemark}
  \footnotetextmetadata{#2}%
}

%%% End of preamble-accessible-graphics.tex


%% ============================================================================
%% Document Info
%% ============================================================================
\title{Accessibility Test Document\\{\large Alt Text for Images via BibLaTeX Metadata}}
\author{Robert Hildebrand}
\date{\today}

%% ============================================================================
%% Document Body
%% ============================================================================
\begin{document}

\maketitle

\begin{abstract}
This document demonstrates the use of accessible graphics commands that automatically
pull alt text descriptions from BibLaTeX metadata files. The alt text is stored in
the \texttt{abstract} field of each bibliography entry and is embedded into the PDF
for screen reader compatibility.
\end{abstract}

\tableofcontents
\newpage

%% ============================================================================
\section{Introduction}
%% ============================================================================

This document tests the accessible graphics functionality. Each image command
automatically:
\begin{enumerate}
  \item Includes the image from the appropriate figures folder
  \item Adds alt text from the \texttt{abstract} field in the \texttt{.bib} file
  \item Displays copyright attribution
  \item Adds a footnote with full metadata (title, URL, author, date)
\end{enumerate}

%% ============================================================================
\section{Inline Graphics (No Figure Environment)}
%% ============================================================================

The \verb|\includegraphicstatic| command includes an image inline (centered, but
not in a float). Here is an example showing a linear programming feasible region:

\includegraphicstatic[width=0.6\linewidth]{LP-feasible-region.png}

The alt text for this image (stored in the bib file) describes the polygonal
feasible region and constraint lines.

%% ============================================================================
\section{Inline Graphics from figures-source}
%% ============================================================================

The \verb|\includegraphicsource| command includes an image from the figures-source folder.
Here is an example showing an undirected graph:

\includegraphicsource[width=0.5\linewidth]{tikz-modeling-sums-continued-01-0564d3a1}

This tests that the alt text and metadata are correctly pulled from the figures-source
metadata file.

%% ============================================================================
\section{No-Cite Versions (No Attribution)}
%% ============================================================================

For graphics you created yourself where attribution isn't needed, use the ``nocite''
versions. These still include alt text for accessibility but omit the copyright line
and footnote metadata.

\subsection{Inline No-Cite (figures-source)}

The \verb|\includegraphicsourcenocite| command:

\includegraphicsourcenocite[width=0.5\linewidth]{tikz-modeling-sums-continued-01-0564d3a1}

Notice there's no copyright line or footnote below the image.

\subsection{Inline No-Cite (figures-static)}

The \verb|\includegraphicstaticnocite| command works the same way for images in the figures-static folder:

\includegraphicstaticnocite[width=0.5\linewidth]{wiki-File-integer-programming.png}

This shows the IP polytope with LP relaxation, with alt text but no attribution.

\subsection{Figure Environment No-Cite}

The \verb|\includefiguresourcenocite| command creates a figure with caption but no attribution:

\includefiguresourcenocite[Undirected graph example without attribution][width=0.5\linewidth][ht]{tikz-modeling-sums-continued-01-0564d3a1}

%% ============================================================================
\section{Figure Environment with Auto-Caption}
%% ============================================================================

The \verb|\includefigurestatic| command creates a proper figure environment with
a caption pulled from the bibliography title field.

% Arguments: [caption][graphics options][float placement]{bibkey}
\includefigurestatic[][width=0.6\linewidth][ht]{IP-feasible-region.png}

You can reference this figure using \verb|\autoref{fig:IP-feasible-region.png}|,
which produces: \autoref{fig:IP-feasible-region.png}.

%% ============================================================================
\section{Figure with Custom Caption}
%% ============================================================================

You can also provide a custom caption as the first optional argument:

% Arguments: [caption][graphics options][float placement]{bibkey}
\includefigurestatic[A support vector machine separating two classes of data][width=0.5\linewidth][ht]{SVM-plot.png}

The alt text still comes from the bib file's abstract field, but the caption
is overridden.

%% ============================================================================
\section{Manual Alt Text Override}
%% ============================================================================

For cases where you need to provide alt text directly (perhaps context-specific),
use the \verb|\includegraphicstaticalt| command:

\includegraphicstaticalt[width=0.5\linewidth]{LP-feasible-region.png}{A simple 2D plot showing the intersection of linear constraints forming a bounded polygon. The optimal point is at the top-right vertex where two constraint lines meet.}

This overrides the alt text from the bib file while keeping all other metadata.

%% ============================================================================
\section{How It Works}
%% ============================================================================

\subsection{The Metadata File}

Alt text is stored in the \texttt{abstract} field of each bib entry:

\begin{verbatim}
@Online{LP-feasible-region.png,
  author = {{Robert Hildebrand [CC BY-SA 4.0]}},
  title = {Linear Programming Feasible Region},
  abstract = {A two-dimensional coordinate system
    showing a shaded polygonal region representing
    the feasible region of a linear program...},
  year = {2023},
  options = {skipbib=true},
}
\end{verbatim}

\subsection{The Document Setup}

The key requirement is the \verb|\DocumentMetadata| command \emph{before}
\verb|\documentclass|:

\begin{verbatim}
% For TeX Live 2025 (January release):
\DocumentMetadata{
  lang=en-US,
  pdfversion=2.0,
  testphase={phase-III}
}
\documentclass{article}

% For later releases (2025-11-01+), use:
% \DocumentMetadata{lang=en-US, pdfversion=2.0, tagging=on}
\end{verbatim}

\subsection{Available Commands}

\begin{table}[ht]
\centering
\begin{tabular}{llll}
\hline
\textbf{Command} & \textbf{Folder} & \textbf{Float} & \textbf{Citation} \\
\hline
\verb|\includegraphicstatic| & figures-static/ & No & Yes \\
\verb|\includegraphicsource| & figures-source/ & No & Yes \\
\verb|\includegraphicstaticnocite| & figures-static/ & No & No \\
\verb|\includegraphicsourcenocite| & figures-source/ & No & No \\
\verb|\includefigurestatic| & figures-static/ & Yes & Yes \\
\verb|\includefiguresource| & figures-source/ & Yes & Yes \\
\verb|\includefigurestaticnocite| & figures-static/ & Yes & No \\
\verb|\includefiguresourcenocite| & figures-source/ & Yes & No \\
\verb|\refincludefigurestatic| & figures-static/ & Yes + autoref & Yes \\
\verb|\refincludefiguresource| & figures-source/ & Yes + autoref & Yes \\
\verb|\includegraphicstaticalt| & figures-static/ & No (manual alt) & Yes \\
\verb|\includegraphicsourcealt| & figures-source/ & No (manual alt) & Yes \\
\hline
\end{tabular}
\caption{Summary of accessible graphics commands}
\end{table}

%% ============================================================================
\section{Verifying Accessibility}
%% ============================================================================

To verify that alt text is correctly embedded:

\begin{enumerate}
  \item Open the PDF in Adobe Acrobat Pro
  \item Go to \textbf{View} $\rightarrow$ \textbf{Navigation Panels} $\rightarrow$ \textbf{Tags}
  \item Expand the tag tree and find \texttt{<Figure>} tags
  \item Right-click and select \textbf{Properties} to see the alt text
  \item Or use \textbf{Accessibility} $\rightarrow$ \textbf{Full Check} to run an accessibility audit
\end{enumerate}

Alternatively, use a screen reader (NVDA, JAWS, VoiceOver) to test that
images are announced with their descriptions.

%% ============================================================================
%% Bibliography (optional - we use skipbib=true for figures)
%% ============================================================================

\printbibliography

\end{document}
